% =========================================================
% Proof Stub: Bolzano--Weierstrass Theorem
% Source: volume-iii/analysis/sequences/notes/sequences/notes-subsequences.tex
% Mode: standard
% =========================================================

\clearpage
\phantomsection
\hypertarget{proof-RA-ABB-P-BW}{}

\subsubsection[Bolzano--Weierstrass Theorem]{Proof --- Bolzano--Weierstrass Theorem}

\bigskip

\noindent
\textbf{Source.}
Abbott, \emph{Understanding Analysis}, Theorem~2.5.5; Tao, \emph{Analysis I}.
See notes \texttt{notes-subsequences.tex}.

\vspace{0.75em}

\noindent
\textbf{Goal.}
Every bounded sequence in $\mathbb{R}$ has a convergent subsequence.

\vspace{0.75em}

\noindent
\textbf{Logical form.}
\[
  (x_n) \text{ bounded}
  \;\Rightarrow\;
  \exists \text{ subsequence } (x_{n_k}) \text{ and } L \in \mathbb{R}
  \text{ with } (x_{n_k}) \to L.
\]

\vspace{0.75em}

\noindent
\textbf{Proof strategy.}
\textbf{Via Monotone Subsequence Theorem:} Every sequence has a monotone
subsequence. A bounded monotone subsequence converges by MCT.

\textbf{Alternative — Nested Intervals:} Since $(x_n)$ is bounded by
$[a, b]$, repeatedly bisect: at each stage at least one half contains
infinitely many terms; choose that half, select a term, iterate. The
constructed subsequence lies in a sequence of nested intervals with lengths
$\to 0$; by NIP these intervals converge to a single point, which is the
limit.

\vspace{0.75em}

\noindent
\textbf{Note.}
The Nested Intervals proof is more elementary but more involved.
The Monotone Subsequence approach is shorter once you have that lemma.

\vspace{0.75em}

\noindent
\textbf{Proof.}
\begin{proof}
\Claim Every bounded sequence in $\mathbb{R}$ has a convergent subsequence.

\Given $(x_n)$ bounded in $\mathbb{R}$.

\Goal Construct a convergent subsequence.

\Strategy Use the Monotone Subsequence Theorem to extract a monotone
subsequence; apply MCT.

\medskip

% --- by Monotone Subsequence Theorem: (x_n) has a monotone subsequence (x_{n_k}) ---

% --- (x_{n_k}) is bounded (as a subsequence of a bounded sequence) ---

% --- (x_{n_k}) is monotone and bounded, so converges by MCT ---

\end{proof}

\begin{remark*}[Proof shape]
Two-step: extract a monotone subsequence (Monotone Subsequence Theorem),
then invoke MCT for convergence.
\end{remark*}

\begin{remark*}[Dependencies]
\begin{itemize}
  \item Monotone Subsequence Theorem: every sequence has a monotone subsequence.
  \item Monotone Convergence Theorem (MCT).
  \item Boundedness is inherited by subsequences.
\end{itemize}
\end{remark*}
